\section{Benchmark Corpus}
\label{app:corpus}

The benchmark comprises the 40 photon and neutron science papers listed below,
ordered by arXiv identifier. Each entry links to the public arXiv record.

\begin{small}
\begin{enumerate}\setlength{\itemsep}{1pt}
  \item \href{https://arxiv.org/abs/1804.07163}{arXiv:1804.07163}: Insights into the interactions among Surfactin, betaines, and PAM: surface tension, small-angle neutron scattering, and small-angle X-ray scattering study
  \item \href{https://arxiv.org/abs/1806.08116}{arXiv:1806.08116}: The structure of near stoichiometric Ge-Ga-Sb-S glasses: a reverse Monte Carlo study
  \item \href{https://arxiv.org/abs/1811.10341}{arXiv:1811.10341}: Local order in binary Ge-Te glasses - an experimental study
  \item \href{https://arxiv.org/abs/1812.03956}{arXiv:1812.03956}: X-ray reflectivity with a twist: quantitative time-resolved X-ray reflectivity using monochromatic synchrotron radiation
  \item \href{https://arxiv.org/abs/1901.08819}{arXiv:1901.08819}: Polarized neutron reflectometry study from iron oxide nanoparticles monolayer
  \item \href{https://arxiv.org/abs/1907.00342}{arXiv:1907.00342}: The X-ray photoemission and Co K -Edge X-ray absorption of Ba2CoWO6
  \item \href{https://arxiv.org/abs/1909.08292}{arXiv:1909.08292}: Neutron diffraction and magnetic properties of Co$_2$Cr$_{1-x}$Ti$_x$Al Heusler alloys
  \item \href{https://arxiv.org/abs/1912.11816}{arXiv:1912.11816}: Experimental observation of third-order effect in magnetic small-angle neutron scattering
  \item \href{https://arxiv.org/abs/2001.00707}{arXiv:2001.00707}: Neutron diffraction and ab initio studies on the fully compensated ferrimagnetic characteristics of Mn2V1-xCoxGa Heusler alloys
  \item \href{https://arxiv.org/abs/2001.02706}{arXiv:2001.02706}: Magnetic correlations in the disordered ferromagnetic alloy Ni-V revealed with small angle neutron scattering
  \item \href{https://arxiv.org/abs/2005.01985}{arXiv:2005.01985}: Mechanism of pressure induced amorphization of SnI4: a combined X-ray diffraction -- X-ray absorption spectroscopy study
  \item \href{https://arxiv.org/abs/2006.12074}{arXiv:2006.12074}: Magnetic excitations of Sr$_3$Ir$_2$O$_7$ observed by inelastic neutron scattering measurement
  \item \href{https://arxiv.org/abs/2007.00972}{arXiv:2007.00972}: Synthesis and Characterization of Strontium Doped Barium Titanates using Neutron Diffraction Technique
  \item \href{https://arxiv.org/abs/2009.01093}{arXiv:2009.01093}: X-ray linear dichroic ptychography
  \item \href{https://arxiv.org/abs/2009.12925}{arXiv:2009.12925}: Dynamical spin susceptibility in La2CuO4 studied by resonant inelastic x-ray scattering
  \item \href{https://arxiv.org/abs/2010.00094}{arXiv:2010.00094}: Local and electronic structure of Sr1-xGdxTiO3 probed by X-ray absorption spectroscopy
  \item \href{https://arxiv.org/abs/2010.02965}{arXiv:2010.02965}: New insights into the mechanism of graphene oxide and radionuclide interaction through vacancy defects
  \item \href{https://arxiv.org/abs/2010.12832}{arXiv:2010.12832}: X-ray tomography investigation of cyclically sheared granular materials
  \item \href{https://arxiv.org/abs/2011.13326}{arXiv:2011.13326}: In-situ investigation of temperature induced agglomeration in non-polar magnetic nanoparticle dispersions by small angle X-ray scattering
  \item \href{https://arxiv.org/abs/2104.10884}{arXiv:2104.10884}: X-ray Absorption Spectroscopy Study of La$_{1-y}$Sr$_{y}$Co$_{1-x}$Nb$_x$O$_3$
  \item \href{https://arxiv.org/abs/2105.04167}{arXiv:2105.04167}: Role of higher-order effects in spin-misalignment small-angle neutron scattering of high-pressure torsion nickel
  \item \href{https://arxiv.org/abs/2202.11658}{arXiv:2202.11658}: The honeycomb and hyperhoneycomb polymorphs of IrI$_3$
  \item \href{https://arxiv.org/abs/2203.07170}{arXiv:2203.07170}: In situ study of zinc peroxide decomposition to zinc oxide by X-ray absorption spectroscopy and reverse Monte-Carlo simulations
  \item \href{https://arxiv.org/abs/2209.06540}{arXiv:2209.06540}: Short range order and topology of binary Ge-S glasses
  \item \href{https://arxiv.org/abs/2211.06006}{arXiv:2211.06006}: Small-angle neutron scattering of long-wavelength magnetic modulations in reduced sample dimensions
  \item \href{https://arxiv.org/abs/2212.04169}{arXiv:2212.04169}: Structural changes in Ti1-xAlxN coatings during turning: A XANES and EXAFS study of worn tools
  \item \href{https://arxiv.org/abs/2302.02010}{arXiv:2302.02010}: Identifying f-electron symmetries of UTe2 with O-edge resonant inelastic X-ray scattering
  \item \href{https://arxiv.org/abs/2304.08002}{arXiv:2304.08002}: Morphology of buried interfaces in ion-assisted magnetron sputter deposited 11B4C-containing Ni/Ti multilayer neutron optics investigated by grazing incidence small angle scattering
  \item \href{https://arxiv.org/abs/2309.13729}{arXiv:2309.13729}: Approaching Ferrite-Based Exchange-Coupled Nanocomposites as Permanent Magnets
  \item \href{https://arxiv.org/abs/2404.06853}{arXiv:2404.06853}: Revealing mechanism of pore defect formation in laser directed energy deposition of aluminum alloy via in-situ synchrotron X-ray imaging
  \item \href{https://arxiv.org/abs/2405.14324}{arXiv:2405.14324}: Magnetic microstructure of nanocrystalline Fe-Nb-B alloys as seen by small-angle neutron and X-ray scattering
  \item \href{https://arxiv.org/abs/2405.15411}{arXiv:2405.15411}: Short range order and topology of Te-rich amorphous Ge-Sb-Te alloys
  \item \href{https://arxiv.org/abs/2407.10282}{arXiv:2407.10282}: High-speed synchrotron X-ray imaging of melt pool dynamics during ultrasonic melt processing of Al6061
  \item \href{https://arxiv.org/abs/2407.20626}{arXiv:2407.20626}: Electronic structure and resonant inelastic x-ray scattering in Ta2NiSe5
  \item \href{https://arxiv.org/abs/2410.00268}{arXiv:2410.00268}: Operando Ultrafast Damage-free Diffraction-Enhanced X-ray Absorption Spectroscopy for Chemical Reactivity of Polymer in Solvents
  \item \href{https://arxiv.org/abs/2501.08122}{arXiv:2501.08122}: Origin of perpendicular magnetic anisotropy in ultra-thin metal films studied by in-situ neutron reflectometry
  \item \href{https://arxiv.org/abs/2503.18475}{arXiv:2503.18475}: Rietveld Refinement and NMR Crystallographic Investigations of Multicomponent Crystals Containing Alkali Metal Chlorides and Urea
  \item \href{https://arxiv.org/abs/2601.16215}{arXiv:2601.16215}: High-resolution neutron diffraction determination of noncollinear antiferromagnetic order in the honeycomb magnetoelectric Fe$_{4}$Nb$_{2}$O$_{9}$
  \item \href{https://arxiv.org/abs/2603.01364}{arXiv:2603.01364}: Probing Coordination Environments in Buried Oxides of Aluminum Josephson Junctions by Resonant X-ray Reflectivity
  \item \href{https://arxiv.org/abs/2605.14792}{arXiv:2605.14792}: Thermal expansion of FeWO$_4$ (Ferberite) and FeWO$_4$:Fe$_2$WO$_6$ (7:1): a comparative X-ray and neutron diffraction study
\end{enumerate}
\end{small}


\section{Schema}
\label{app:schema}

The extraction target is defined by a two-layer schema: a generic layer and a PaN
layer. Every field carries a description and a type; top-level fields also declare
a cardinality, and most fields carry an \texttt{instructions} array defining inclusion,
normalization, or entity-granularity decisions. Repeated entities define an
\texttt{item\_schema} for their subfields. Figure~\ref{fig:schema-excerpt} shows
two representative excerpts: the paper-level \texttt{title} field and the
\texttt{samples} entity. The complete set of \emph{scored} fields is enumerated
in the field-level results (Tables~\ref{tab:appendix-iaa} and~\ref{tab:appendix-zero-shot});
the full per-field definitions appear only in the released schema. The no-instruction
variants used in the instruction ablation (Section~\ref{sec:instruction-ablation}) simply
remove the \texttt{instructions} arrays.

\begin{figure*}[t]
\textbf{Paper-level field} (\texttt{bibliographic.title}):
\begin{lstlisting}[style=codeblock]
{
  "title": {
    "description": "Main title of the scientific document itself.",
    "type": "string|null",
    "cardinality": {
      "min": 0,
      "max": 1
    },
    "instructions": [
      "Copy the document title verbatim as printed on the first page or running header. Preserve casing and punctuation."
    ]
  }
}
\end{lstlisting}
\textbf{Entity with sub-schema} (\texttt{samples}, abridged):
\begin{lstlisting}[style=codeblock]
{
  "samples": {
    "description": "Samples studied through experimental measurements reported or reanalyzed in the document. A sample entry can represent a single specimen or a measured group when only the group is identified by the document.",
    "type": "list[sample]",
    "cardinality": {
      "min": 0,
      "max": 50
    },
    "instructions": [
      "Follow this decision routine to determine sample entries.",
      "..."
    ],
    "item_schema": {
      "id": {
        "type": "string",
        "description": "Stable local sample identifier.",
        "instructions": [
          "Assign IDs in sequence: sample_1, sample_2, sample_3, etc.",
          "..."
        ]
      },
      "sample_name": {
        "type": "string",
        "description": "Name, label, or concise identifying phrase for the sample.",
        "instructions": [
          "Follow this decision routine to extract sample_name.",
          "..."
        ]
      },
      "...": "5 further subfields: sample_description, chemical_composition, sample_state, sample_mass, supplier_or_creator"
    }
  }
}
\end{lstlisting}
  \caption{Two excerpts from the released two-layer schema: a paper-level
  field shown in full, and an entity collection with
  its \texttt{item\_schema} of subfields. Each field carries a
  \texttt{description} and \texttt{type}; top-level fields also declare
  a \texttt{cardinality}, and most fields carry an \texttt{instructions}
  array (removed by the no-instruction ablation of
  Section~\ref{sec:instruction-ablation}). The \texttt{samples} entry is
  abridged: long instruction routines and the remaining subfields are
  elided (marked \texttt{...}), and the complete schema with all
  per-field definitions ships with the benchmark.}
  \label{fig:schema-excerpt}
\end{figure*}


\section{Prompt Template}
\label{app:prompt}

All conditions use a single text-to-JSON prompt. The document is converted to
Markdown-like text and inserted together with the schema and an output
skeleton; the model returns one JSON object. Figure~\ref{fig:prompt-template}
shows the zero-shot template.

\begin{figure*}[t]
\begin{lstlisting}[style=codeblock,basicstyle=\ttfamily\scriptsize]
## Task
You are an expert system for extracting metadata from scientific documents.

Extract metadata from the document according to the provided schema.

## Schema

{{SCHEMA_JSON}}

## Rules

- Use only information explicitly stated in the document; do not infer or guess missing information.
- Extract metadata only for the document itself; ignore cited references and other documents mentioned in the text.
- Preserve values exactly as written unless the schema explicitly requires normalization.
- Follow the schema structure strictly.

- For fields with type `string|null`:
  - use a string when the value is present
  - use `null` when the value is not found

- For fields with list types:
  - use a list of extracted items when present
  - use `[]` when no items are found

- For fields with type `boolean`:
  - use `true` when clearly supported by the document
  - use `false` when the document explicitly states the contrary

- For nested objects and list items:
  - follow the schema structure exactly
  - fill subfields according to their declared types

## Output Format

Return a single valid JSON object matching the skeleton below exactly.
- Use double-quoted object keys and double-quoted string values.
- Do not wrap the JSON in code fences (```), markdown, or any other formatting.
- Do not add fields not defined in the skeleton.
- Do not omit fields defined in the skeleton.
- Do not return explanations, comments, or natural-language text outside the JSON.

{{OUTPUT_SKELETON}}

## Document

{{DOCUMENT_TEXT}}
\end{lstlisting}
  \caption{Zero-shot extraction prompt template. \texttt{\{\{SCHEMA\_JSON\}\}},
  \texttt{\{\{OUTPUT\_SKELETON\}\}}, and \texttt{\{\{DOCUMENT\_TEXT\}\}} are
  replaced at request time. The few-shot template additionally inserts an
  \texttt{Examples} block of complete worked extractions before the
  document.}
  \label{fig:prompt-template}
\end{figure*}


\section{Reproducibility Details}
\label{app:reproducibility}

All 660 extractions included in the reported evaluation use one text-to-JSON
pipeline under the fixed settings in Table~\ref{tab:setup}. Each of the 22
conditions contains exactly one extraction per target paper, so all reported
scores are single-run point estimates without seed averaging. The seven model
snapshots are pinned to the versions in
Table~\ref{tab:experimental-matrix}.

\begin{table}[t]
  \centering
  \small
  \begin{tabular}{@{}ll@{}}
    \toprule
    Setting & Value \\
    \midrule
    \multicolumn{2}{@{}l}{\textit{Document input}} \\
    PDF-to-text & \texttt{pymupdf4llm} 1.27.2.2 \\
    \addlinespace
    \multicolumn{2}{@{}l}{\textit{Decoding}} \\
    Temperature & 0 \\
    Reasoning / thinking & Disabled \\
    Max output tokens & 12{,}288 \\
    Runs per condition & 1 \\
    \addlinespace
    \multicolumn{2}{@{}l}{\textit{Local serving (Gemma 4, Qwen3.5)}} \\
    Engine & vLLM 0.24.0 \\
    GPU & NVIDIA A100 80GB \\
    \addlinespace
    \multicolumn{2}{@{}l}{\textit{Hosted serving (GPT-5.5, GPT-4o, Mistral)}} \\
    Access & OpenAI-compatible gateway \\
    \bottomrule
  \end{tabular}
  \caption{Fixed pipeline settings applied to all 22 conditions.
  Per-model serving mode and the pinned model
  snapshots are listed in Table~\ref{tab:experimental-matrix}.}
  \label{tab:setup}
\end{table}


The schema, prompt template, reconciled gold records, scoring plan, alignment
helper, and evaluation code are released with the benchmark.


\section{Complete Field-Level Results}
\label{app:complete-results}

This appendix reports the complete field-level evaluation results. Section
headings have separate rows for content recovery and relative order; all other
fields appear once. The metric used for each row is shown in the tables.
Fields with no or very few non-null gold values are flagged
($^{\dagger}$/$^{\ddagger}$). Their scores are reported for completeness but
are not interpreted.

The three F1 variants use the same formula but differ in how matches are
determined: set F1 uses exact matching after normalization, entity F1 uses the
one-to-one alignment described in Section~\ref{sec:entity-alignment}, and
reference F1 is computed after predicted output-local identifiers are remapped
through that alignment.

Table~\ref{tab:appendix-iaa} reports the complete pairwise annotation agreement
results for the 30-paper evaluation set.

\begin{table*}[t]
  \centering
  \footnotesize
  \setlength{\tabcolsep}{6pt}
  \begin{tabular}{@{}p{0.34\textwidth}p{0.18\textwidth}r@{}}
    \toprule
     & Metric & Agreement \\
    \midrule
    \multicolumn{3}{@{}l}{\textit{Paper-level fields}} \\
    Title & ROUGE-L & 0.957 \\
    Subtitle$^{\dagger}$ & ROUGE-L & 1.000 \\
    Abstract & ROUGE-L & 0.996 \\
    DOI$^{\dagger}$ & Normalized exact & 1.000 \\
    arXiv ID & Normalized exact & 1.000 \\
    URL$^{\dagger}$ & Normalized exact & 1.000 \\
    Publication date & Exact & 0.900 \\
    Venue & Normalized exact & 1.000 \\
    Publisher$^{\dagger}$ & Normalized exact & 1.000 \\
    Language & Exact & 1.000 \\
    Keywords & Set F1 & 0.976 \\
    Section-heading content & Set F1 & 0.969 \\
    Section-heading order & Kendall $\tau$ & 1.000 \\
    Has references & Exact & 1.000 \\
    Has figures & Exact & 1.000 \\
    Has tables & Exact & 1.000 \\
    \addlinespace
    \midrule
    \multicolumn{3}{@{}l}{\textit{Affiliations}} \\
    Affiliation entities & Entity F1 & 1.000 \\
    Affiliation label & Normalized exact & 0.962 \\
    Affiliation text & ROUGE-L & 0.997 \\
    \addlinespace
    \midrule
    \multicolumn{3}{@{}l}{\textit{Authors}} \\
    Author entities & Entity F1 & 1.000 \\
    Author full name & Normalized exact & 0.996 \\
    Author ORCID$^{\dagger}$ & Normalized exact & 1.000 \\
    Author email & Normalized exact & 1.000 \\
    Corresponding-author flag & Exact & 0.987 \\
    Author--affiliation refs & Reference F1 & 1.000 \\
    \addlinespace
    \midrule
    \multicolumn{3}{@{}l}{\textit{Funding}} \\
    Funding entities & Entity F1 & 0.889 \\
    Funder & ROUGE-L & 0.837 \\
    Grant ID & Normalized exact & 0.900 \\
    \addlinespace
    \midrule
    \multicolumn{3}{@{}l}{\textit{Facilities}} \\
    Facility entities & Entity F1 & 0.753 \\
    Facility name & ROUGE-L & 0.776 \\
    \addlinespace
    \midrule
    \multicolumn{3}{@{}l}{\textit{Samples}} \\
    Sample entities & Entity F1 & 0.845 \\
    Sample name & ROUGE-L & 0.740 \\
    Chemical composition & Normalized exact & 0.817 \\
    Sample state & Normalized exact & 0.718 \\
    Sample mass$^{\ddagger}$ & Normalized exact & 0.986 \\
    Sample description & ROUGE-L & 0.180 \\
    Supplier / creator & ROUGE-L & 0.845 \\
    \addlinespace
    \midrule
    \multicolumn{3}{@{}l}{\textit{Datasets}} \\
    Dataset entities & Entity F1 & 0.793 \\
    Dataset name & ROUGE-L & 0.693 \\
    Measurement technique & ROUGE-L & 0.864 \\
    Dataset description & ROUGE-L & 0.361 \\
    Measurement mode & Set F1 & 0.700 \\
    Dataset purpose & Exact & 0.957 \\
    Dataset--sample refs & Reference F1 & 0.855 \\
    Dataset--facility refs & Reference F1 & 0.814 \\
    \bottomrule
  \end{tabular}
  \caption{Complete pairwise annotation agreement on the 30
  evaluation papers. Attribute and reference rows are
  conditional on aligned parents. A $^{\dagger}$ marks a field with
  no non-null support and a $^{\ddagger}$ one with at most 10\%
  non-null support.}
  \label{tab:appendix-iaa}
\end{table*}


Table~\ref{tab:appendix-zero-shot} reports the complete field-level results for
all seven models under the full-instruction zero-shot setting on the 30-paper
evaluation set.

Three paper-level fields marked $^{\S}$ are affected by the PDF-to-text
conversion: arXiv ID, publication date, and venue. On affected papers, the
relevant evidence appears only in the rotated first-page arXiv stamp, which
\texttt{pymupdf4llm} omits from the model input. Most notably, the arXiv-ID
score of 0.533 reflects null--null agreement on 16 papers and no successful
recovery on the 14 papers with non-null gold values. Scores for these fields
therefore partly reflect unavailable input evidence rather than model
extraction alone.

\begin{table*}[t]
  \centering
  \fontsize{8.5pt}{9.5pt}\selectfont
  \setlength{\tabcolsep}{2pt}
  \begin{tabular}{@{}p{0.24\textwidth}p{0.13\textwidth}rrrrrrr@{}}
    \toprule
     & Metric & GPT-5.5 & GPT-4o & Mistral & Gemma & Qwen 27B & Qwen 9B & Qwen 4B \\
    \midrule
    \multicolumn{9}{@{}l}{\textit{Paper-level fields}} \\
    Title & ROUGE-L & 0.994 & 0.992 & 0.995 & 0.990 & 0.990 & 0.992 & 0.993 \\
    Subtitle$^{\dagger}$ & ROUGE-L & 0.967 & 0.967 & 1.000 & 1.000 & 1.000 & 1.000 & 1.000 \\
    Abstract & ROUGE-L & 0.976 & 0.987 & 0.982 & 0.973 & 0.974 & 0.980 & 0.977 \\
    DOI$^{\dagger}$ & Normalized exact & 1.000 & 0.967 & 1.000 & 0.967 & 1.000 & 0.967 & 1.000 \\
    arXiv ID$^{\S}$ & Normalized exact & 0.533 & 0.533 & 0.533 & 0.533 & 0.533 & 0.533 & 0.533 \\
    URL$^{\dagger}$ & Normalized exact & 1.000 & 1.000 & 1.000 & 1.000 & 1.000 & 1.000 & 1.000 \\
    Publication date$^{\S}$ & Exact & 0.800 & 0.767 & 0.667 & 0.800 & 0.800 & 0.767 & 0.800 \\
    Venue$^{\S}$ & Normalized exact & 0.567 & 0.567 & 0.600 & 0.567 & 0.567 & 0.567 & 0.667 \\
    Publisher$^{\dagger}$ & Normalized exact & 1.000 & 0.933 & 0.933 & 0.967 & 0.967 & 0.933 & 0.967 \\
    Language & Exact & 1.000 & 1.000 & 1.000 & 1.000 & 1.000 & 0.967 & 1.000 \\
    Keywords & Set F1 & 1.000 & 1.000 & 1.000 & 1.000 & 1.000 & 0.976 & 0.909 \\
    Section-heading content & Set F1 & 0.945 & 0.763 & 0.660 & 0.949 & 0.969 & 0.635 & 0.710 \\
    Section-heading order & Kendall $\tau$ & 1.000 & 1.000 & 1.000 & 1.000 & 1.000 & 1.000 & 1.000 \\
    Has references & Exact & 1.000 & 1.000 & 1.000 & 1.000 & 1.000 & 1.000 & 1.000 \\
    Has figures & Exact & 1.000 & 1.000 & 1.000 & 1.000 & 1.000 & 1.000 & 1.000 \\
    Has tables & Exact & 1.000 & 1.000 & 0.967 & 1.000 & 1.000 & 1.000 & 1.000 \\
    \addlinespace
    \midrule
    \multicolumn{9}{@{}l}{\textit{Affiliations}} \\
    Affiliation entities & Entity F1 & 1.000 & 1.000 & 1.000 & 1.000 & 1.000 & 1.000 & 1.000 \\
    Affiliation label & Normalized exact & 1.000 & 0.977 & 1.000 & 1.000 & 1.000 & 1.000 & 0.970 \\
    Affiliation text & ROUGE-L & 0.999 & 0.997 & 0.999 & 0.999 & 0.999 & 0.999 & 0.999 \\
    \addlinespace
    \midrule
    \multicolumn{9}{@{}l}{\textit{Authors}} \\
    Author entities & Entity F1 & 1.000 & 1.000 & 1.000 & 1.000 & 1.000 & 1.000 & 1.000 \\
    Author full name & Normalized exact & 1.000 & 0.991 & 1.000 & 1.000 & 1.000 & 1.000 & 1.000 \\
    Author ORCID$^{\dagger}$ & Normalized exact & 1.000 & 1.000 & 1.000 & 1.000 & 1.000 & 1.000 & 1.000 \\
    Author email & Normalized exact & 0.996 & 0.996 & 0.996 & 0.996 & 0.996 & 0.991 & 0.978 \\
    Corresponding-author flag & Exact & 0.991 & 0.991 & 1.000 & 0.982 & 1.000 & 0.982 & 0.965 \\
    Author--affiliation refs & Reference F1 & 1.000 & 1.000 & 1.000 & 1.000 & 1.000 & 0.994 & 0.996 \\
    \addlinespace
    \midrule
    \multicolumn{9}{@{}l}{\textit{Funding}} \\
    Funding entities & Entity F1 & 0.954 & 0.925 & 0.912 & 0.945 & 0.953 & 0.901 & 0.810 \\
    Funder & ROUGE-L & 0.872 & 0.833 & 0.822 & 0.892 & 0.885 & 0.829 & 0.818 \\
    Grant ID & Normalized exact & 0.694 & 0.855 & 0.790 & 0.850 & 0.770 & 0.875 & 0.878 \\
    \addlinespace
    \midrule
    \multicolumn{9}{@{}l}{\textit{Facilities}} \\
    Facility entities & Entity F1 & 0.545 & 0.892 & 0.735 & 0.860 & 0.769 & 0.783 & 0.744 \\
    Facility name & ROUGE-L & 0.820 & 0.747 & 0.806 & 0.816 & 0.855 & 0.767 & 0.812 \\
    \addlinespace
    \midrule
    \multicolumn{9}{@{}l}{\textit{Samples}} \\
    Sample entities & Entity F1 & 0.800 & 0.777 & 0.865 & 0.864 & 0.837 & 0.783 & 0.739 \\
    Sample name & ROUGE-L & 0.788 & 0.840 & 0.734 & 0.854 & 0.884 & 0.772 & 0.783 \\
    Chemical composition & Normalized exact & 0.750 & 0.689 & 0.688 & 0.645 & 0.611 & 0.615 & 0.655 \\
    Sample state & Normalized exact & 0.486 & 0.443 & 0.532 & 0.658 & 0.667 & 0.569 & 0.552 \\
    Sample mass$^{\ddagger}$ & Normalized exact & 0.944 & 0.852 & 0.987 & 0.947 & 0.986 & 0.954 & 0.966 \\
    Sample description & ROUGE-L & 0.403 & 0.240 & 0.280 & 0.444 & 0.306 & 0.326 & 0.261 \\
    Supplier / creator & ROUGE-L & 0.846 & 0.790 & 0.675 & 0.781 & 0.847 & 0.747 & 0.738 \\
    \addlinespace
    \midrule
    \multicolumn{9}{@{}l}{\textit{Datasets}} \\
    Dataset entities & Entity F1 & 0.832 & 0.609 & 0.656 & 0.778 & 0.814 & 0.732 & 0.636 \\
    Dataset name & ROUGE-L & 0.717 & 0.576 & 0.723 & 0.722 & 0.732 & 0.702 & 0.510 \\
    Measurement technique & ROUGE-L & 0.747 & 0.903 & 0.871 & 0.886 & 0.950 & 0.863 & 0.839 \\
    Dataset description & ROUGE-L & 0.444 & 0.277 & 0.371 & 0.333 & 0.390 & 0.333 & 0.412 \\
    Measurement mode & Set F1 & 0.752 & 0.475 & 0.463 & 0.590 & 0.681 & 0.333 & 0.495 \\
    Dataset purpose & Exact & 1.000 & 1.000 & 0.984 & 1.000 & 0.958 & 1.000 & 1.000 \\
    Dataset--sample refs & Reference F1 & 0.743 & 0.812 & 0.849 & 0.831 & 0.832 & 0.693 & 0.802 \\
    Dataset--facility refs & Reference F1 & 0.500 & 0.909 & 0.767 & 0.851 & 0.771 & 0.809 & 0.746 \\
    \bottomrule
  \end{tabular}
  \caption{Complete full-instruction zero-shot results on the 30 evaluation
  papers. Paper-level scalar and order scores are averaged
  over papers; scalar entity-attribute scores are averaged over aligned
  pairs; count-based set, entity, and reference F1 scores are
  micro-aggregated across the corpus. Attribute and reference scores are
  conditional
  on aligned parents. A $^{\dagger}$ marks a field
  with no gold non-null values and a $^{\ddagger}$ one with at most
  10\% gold non-null support. A $^{\S}$ marks fields whose gold evidence is
  discarded during PDF-to-text conversion on some papers, as
  detailed in Appendix~\ref{app:complete-results}.}
  \label{tab:appendix-zero-shot}
\end{table*}

